\scanpage{34}
% archive foliation: 33

entendu) les sommets s'identifient aux sections partielles sur des parties \inserted{de $J$} réduites à un élément — la relation d'incidence est celle de prolongement des sections, ou encore l'inclusion des parties de $\Delta$.

Le groupe $G=GL(M)$ opère sur $\Delta$ par autom. d'octaèdre, on voit que $Z=\{\pm1\}$ opère trivialement, \deleted{\unsure{d'où}} et $=$ que c'est là le noyau de l'opération, on a donc un hom injectif
\[
\diagram{\begin{tikzcd}
GL(M)'\arrow[r]\arrow[d,equal]&\mathrm{Aut}_{\mathrm{oct}}(\Delta)\\
GL(M)/\{\pm1\}&
\end{tikzcd}}
\]
\margin{\uline{NB} $GL(M)'\xrightarrow{\sim}GP(M)$ (\unsure{groupe projectif de degré $1$|groupe projectif de dim $1$})}

qui est un isomorphisme, car les deux groupes sont de \deleted{\unsure{un}} même ordre $24\times2=48$, donc on trouve
\[
(33)\qquad
\diagram{\begin{tikzcd}
GL(M)'\arrow[r,"\sim"]\arrow[d,equal]&
\mathrm{Aut}_{\mathrm{oct}}(\Delta)\simeq\mathrm{Aut}^+(\Delta)\times
\underbrace{\{1,\underset{\substack{\downarrow\\\text{antipodisme}}}{\underline a_\Delta}\}}
_{\text{centre de }\mathrm{Aut}_{\mathrm{oct}}(\Delta)}\\
GL(M)/(Z=\{\pm1\})\arrow[d,equal]&\\
GP(M)&
\end{tikzcd}}
\]

Le centre de $GL(M)'$ \deleted{\unsure{correspond}}\inserted{n'est autre} d'ailleurs que le sous-groupe $\mathbb{F}_2^{\underline\omega}/\mathbb{F}_2\cdot\mathrm{id}_{M_0}$ de $GL(M)'$ (car le centre de $GL(M_0)\simeq\unsure{\mathcal{G}_J|\mathfrak S_J}$ est nul, donc le centre de $GL(M)'$ est contenu dans $\underline{\mathrm{End}}(M_0)/(Z=\{0_{M_0},1_{M_0}\})\simeq M_0\times\mathbb{F}_2^{\underline\omega}/\mathbb{F}_2$, et s'identifie au sous-groupe des invariants sous $\unsure{\mathcal{G}_J|\mathfrak S_J}\simeq GL_{\mathbb{F}_2}(M_0)$ dans ce groupe, c'est $\mathbb{F}_2^{\underline\omega}/\mathbb{F}_2$.

La donnée d'un octaèdre $\Delta$ ayant $J$ comme ens des diagonales équivaut à celle d'un torseur
